\documentclass[11pt, twocolumn, landscape, a4paper]{article}

\usepackage[margin=1in]{geometry}
\usepackage[utf8]{inputenc}
\usepackage[spanish]{babel}
\usepackage{amssymb, amsmath, amsbsy}
\usepackage{verbatim}
\usepackage{mathpazo}
\usepackage[pdftex]{hyperref}
\hypersetup{ pdftitle = {ejemplos de derivadas },pdfauthor = {Luis Carrillo Gutiérrez} }

\renewcommand{\familydefault}{\sfdefault}

\begin{document}
\paragraph*{Aplicación de las derivadas}

\begin{itemize}
\item{Sea la función :
$$
f(x) = \begin{cases}
ax^2 + bx + c \qquad  \text{si}\;x \leq 0 \\
\hfill x \cdot \ln x \qquad  \text{si}\;x > 0
\end{cases}
$$

Determinar $a$, $b$ y $c$ para que sea continúa, tenga un máximo en $x = -1$ y la tangente en $x = -2$ sea paralela a la recta $y = 2x$.
}
\end{itemize}
\end{document}